\subsection{Odd-degree gadgets from known powers}
\begin{algorithm}[H]
    \caption{Odd-degree gadget $Q_{2^{l+1}k+(2^l-1)}$ from the known powers
        $(H_2,\ldots,H_{2^l})$, for $k\ge1$ and $l\ge2$.}
    \label{alg:constr-Q-odd}
      \begin{algorithmic}
          \State $\hat{H}_{2^l} = H_{2^l} + Q_{2^{l - 1} - 1}[\alpha_{2^{l + 1} k - 1}, \ldots, \alpha_{2^{l + 1}k + 2^{l - 1} - 3}](x, H_2, \ldots, H_{2^{l - 2}})$
          \State $S^{(1)}_{2^l} = T^{(1)}_{2 k, 2^{l}}[\alpha_{2^{l} - 2}, \ldots, \alpha_{2^{l + 1}k - 3}](x, H_2, \ldots, \hat{H}_{2^l})$
          \State $S^{(2)}_{2^l} = T^{(2)}_{2 k, 2^{l}}[\alpha_{2^{l} - 2}, \ldots, \alpha_{2^{l + 1}k - 3}](x, H_2, \ldots, \hat{H}_{2^l}, \hat{H}_{2^l} + \alpha_{2^{l + 1}k - 2})$
          \State $Q_{2^{l + 1}k + (2^l - 1)}[\alpha_0, \ldots, \alpha_{2^{l + 1}k + 2^l - 2}](x, H_2, \ldots, H_{2^l})$
          \Statex \hspace{\algorithmicindent}$= A_{2^{l-1}}[\alpha_0, \ldots, \alpha_{2^l - 3}, \alpha_{2^{l + 1}k + 2^{l - 1} - 2}, \ldots, \alpha_{2^{l + 1}k + 2^l - 2}](S^{(1)}_{2^l}, S^{(2)}_{2^l}, (x, H_2, \ldots, H_{2^{l-1}}))$
    \end{algorithmic}
\end{algorithm}

\begin{lemma}\label{lem:Q-odd-degree-with-powers}
    Let $k\ge 1$ and $l\ge 2$, put
    $d=2^{l+1}k+(2^l-1)$, and assume that $\mathbb F$ is $d$-admissible.
    Then \Cref{alg:constr-Q-odd} outputs a polynomial $Q_{2^{l+1}k+(2^l-1)}$ that is decodable given $(H_2,\ldots,H_{2^l})$.
\end{lemma}
\begin{proof}
    Let $n\coloneqq\deg S^{(1)}_{2^l}=\deg S^{(2)}_{2^l}=2^{l+1}k$.
    Since $n<d$, $d$-admissibility implies $n$-admissibility.  The internal
    call has parameters $(M,D)=(2k,2^l)$, so its size is exactly $MD=n$;
    moreover $M$ is even, and therefore the exceptional shared odd base is
    irrelevant.  Consequently \Cref{lem:causal-perturbed-T} applies to this
    call.
    The internal $T_{2k,2^l}$ call that defines $(S^{(1)}_{2^l},S^{(2)}_{2^l})$ uses the perturbed power
    $\hat H_{2^l}=H_{2^l}+Q_{2^{l-1}-1}$ (with $\deg Q_{2^{l-1}-1}=2^{l-1}-1$) and its scalar shift
    $\tilde H_{2^l}=\hat H_{2^l}+\alpha_{2^{l+1}k-2}$.
    Apply \Cref{lem:causal-perturbed-T} to this internal call, with
    $D=2^l$, $M=2k$, and $r=2^{l-1}$.  It gives, without treating the
    perturbed power as an already-known constant, that
    \[
      (S^{(1)}_{2^l},S^{(2)}_{2^l})
      \text{ is compatible on }G=\rng{n-r}
    \]
    given $(H_2,\ldots,H_{2^l})$.  The same lemma gives a cutoff-respecting
    recovery of $\hat H_{2^l}$ and $\tilde H_{2^l}$ from the coefficients of
    $xS^{(1)}_{2^l}+S^{(2)}_{2^l}$ and the given powers.
    Thus the internal call is handled directly by
    \Cref{lem:causal-perturbed-T}; no compatibility assertion is imported
    from the recursive program without its cutoff proof.

      Applying \Cref{lem:fill-correctness} at level $l-1$ to the outer $A_{2^{l-1}}$ call shows that the parameters appearing explicitly in that call are extractable from $Q_{2^{l+1}k+(2^l-1)}$ given $(H_2,\ldots,H_{2^l})$, and that $(S^{(1)}_{2^l},S^{(2)}_{2^l})$ is derivable from this given power data.

    It remains to extract the parameters inside the definitions of
    $\hat H_{2^l}$ and the internal $T$ call.  Since the outer fill decoder
    derives $(S^{(1)}_{2^l},S^{(2)}_{2^l})$, it also derives their combined
    polynomial.  The recovery part of \Cref{lem:causal-perturbed-T} therefore
    gives $\hat H_{2^l}$ and $\tilde H_{2^l}$ directly.  Their difference
    extracts the scalar
    $\alpha_{2^{l+1}k-2}=\tilde H_{2^l}-\hat H_{2^l}$.

    Next, from the derived $\hat H_{2^l}$ we obtain the polynomial
    \[
      Q_{2^{l-1}-1}=\hat H_{2^l}-H_{2^l}.
    \]
    This is exactly the $Q_{2^{l-1}-1}$ instance in the definition of $\hat H_{2^l}$, so decodability of $Q_{2^{l-1}-1}$ given $(H_2,\ldots,H_{2^{l-2}})$ extracts the corresponding parameter block.

    Finally, we isolate the combined remainder polynomial for the $T_{2k,2^l}$ call:
    \[
      xR^{(1)}_{2k,2^l}+R^{(2)}_{2k,2^l}
        = \bigl(xS^{(1)}_{2^l}+S^{(2)}_{2^l}\bigr) - \bigl(x\hat H_{2^l}^{2k}+\tilde H_{2^l}^{2k}\bigr).
    \]
    By \Cref{lem:Rk2l}(3), the remaining block of parameters inside $S^{(1)}_{2^l}$ and $S^{(2)}_{2^l}$ is extractable from this polynomial given $(H_2,\ldots,H_{2^{l-1}},\hat H_{2^l},\tilde H_{2^l})$, and all these polynomials are derivable from the data we already have.
    Applying \Cref{lem:discharge-side-information} substitutes those derived
    powers into the conditional remainder decoder.
    Therefore all parameters are extractable from $Q_{2^{l+1}k+(2^l-1)}$ given $(H_2,\ldots,H_{2^l})$, proving that the output is decodable.
\end{proof}

\subsection{Induction steps}

Both the $8k+3$ and $8k+7$ induction steps require auxiliary odd-degree
gadgets that are decodable given only $(H_2,H_4)$.  The special case $n=31$
also needs such a degree-$15$ gadget.  We define the explicit family
$\bar Q_{8k+7}$ for precisely these calls when the available power list
does not contain the next power.  In the induction steps below this choice is
abbreviated by the selected-gadget notation $\mathcal Q_d$.
\Cref{fig:odd-gadgets} summarizes both gadget families and their decoders.

\begin{figure}[H]
  \centering
  \resizebox{\textwidth}{!}{% Ordinary and barred odd-degree gadgets, with the barred decoder blocks.
\begin{tikzpicture}[x=1cm,y=1cm,>=Latex]
  \node[fp panel title] at (0,5.05) {(a) Two ways to obtain an odd auxiliary gadget};
  \node[fp known, minimum width=28mm] (powers) at (1.45,4.18)
    {given powers\\$H_2,\ldots,H_D$};
  \node[fp active, minimum width=32mm] (perturb) at (4.55,4.18)
    {perturb $H_D$ by\\a smaller $Q$-gadget};
  \node[fp recursive, minimum width=25mm] (teven) at (7.65,4.18)
    {even call\\$T_{2k,D}$};
  \node[fp second, minimum width=24mm] (fill) at (10.35,4.18)
    {outer fill\\$A_{D/2}$};
  \node[fp active, minimum width=16mm] (qout) at (12.65,4.18)
    {$Q_d$};
  \draw[fp flow] (powers) -- (perturb);
  \draw[fp flow] (perturb) -- (teven);
  \draw[fp flow] (teven) -- (fill);
  \draw[fp flow] (fill) -- (qout);

  \node[fp known, minimum width=25mm] (h24) at (1.45,3.05)
    {only $H_2,H_4$\\are given};
  \node[fp active, minimum width=30mm] (h8) at (4.55,3.05)
    {construct $H_8$\\with one product};
  \node[fp recursive, minimum width=25mm] (t8) at (7.65,3.05)
    {$T_{k,8}$\\($k=1$ is direct)};
  \node[fp second, minimum width=24mm] (a4) at (10.35,3.05)
    {fixed outer\\fill $A_4$};
  \node[fp active, minimum width=16mm] (barq) at (12.65,3.05)
    {$\bar Q$};
  \draw[fp flow] (h24) -- (h8);
  \draw[fp flow] (h8) -- (t8);
  \draw[fp flow] (t8) -- (a4);
  \draw[fp flow] (a4) -- (barq);

  \node[fp panel title] at (0,2.4) {(b) Top-down decoder of the barred gadget};
  \draw[fp decode] (0,1.75) -- node[fp label, above] {descending coefficient degree}
    (13.55,1.75);
  \node[fp active, minimum width=24mm, anchor=west] at (0,1.08)
    {$b_0,b_1,b_2$\\three unit pivots};
  \node[fp pivot, minimum width=34mm, anchor=west] at (2.4,1.08)
    {$(b_3,b_4,u,v)$\\explicit $4\times4$ block};
  \node[fp pivot, minimum width=19mm, anchor=west] at (5.8,1.08)
    {$w,\rho$\\slopes $k$};
  \node[fp recursive, minimum width=27mm, anchor=west] at (7.7,1.08)
    {internal $R_{k,8}$\\shifted pivots};
  \node[fp second, minimum width=31mm, anchor=west] at (10.4,1.08)
    {remaining $A_4$ slots\\unit pivots};
  \node[fp label, anchor=west] at (2.45,0.42)
    {$\det M=-k^2$ (so $-1$ for $\bar Q_{15}$, the case $k=1$)};
  \node[fp label, anchor=west, text=fpBlue!80!black] at (0,0.02)
    {The matrix inverse is displayed explicitly; no data-dependent division or Jacobian argument is used.};
\end{tikzpicture}
}
  \caption{Ordinary and barred odd-degree gadgets.  When the full power
  list is available, a smaller $Q$ perturbs the top power, an even $T$ call
  supplies a compatible pair, and a fill call closes the gadget.  When only
  $(H_2,H_4)$ is available, the barred construction first synthesizes
  $H_8$ and uses a fixed $A_4$ shell.  Its decoder consists of three unit
  pivots, one explicit constant-determinant block solve, the two power
  shifts, the internal remainder decoder when that block is nonempty, and
  the low fill slots.}
  \label{fig:odd-gadgets}
\end{figure}

\begin{lemma}[A decodable $\bar Q_{8k+7}$ given $(H_2,H_4)$]\label{lem:barQ8k+7}
    Let $k\ge 1$, assume that $\mathbb F$ is $(8k+7)$-admissible, and let
    $H_2,H_4\in\mathbb{F}[x]$ be monic with $\deg H_2=2$ and $\deg H_4=4$.

    Define
    \[
      H_8 \coloneqq \bigl(H_4+(x+\alpha_1)\bigr)\bigl(H_4+(H_2+\alpha_2)\bigr)+\alpha_0,
      \qquad
      \tilde H_8 \coloneqq H_8+\alpha_3.
    \]
    When $k=1$, the interval $\alpha_4,\ldots,\alpha_{8k-5}$ below is empty,
    and the $T$ base is $T^{(1)}_{1,8}=H_8$,
    $T^{(2)}_{1,8}=\widetilde H_8$.
    Let
    \[
      \begin{aligned}
      S^{(1)}&\coloneqq
        T^{(1)}_{k,8}[\alpha_4,\ldots,\alpha_{8k-5}]
          (x,H_2,H_4,H_8),\\
      S^{(2)}&\coloneqq
        T^{(2)}_{k,8}[\alpha_4,\ldots,\alpha_{8k-5}]
          (x,H_2,H_4,H_8,\widetilde H_8),
      \end{aligned}
    \]
    and define
    \[
      \begin{aligned}
      &\bar Q_{8k+7}[\alpha_0,\ldots,\alpha_{8k+6}](x,H_2,H_4)\\
      &\qquad\coloneqq
      A_{4}[\alpha_{8k-4},\ldots,\alpha_{8k+1},
            \alpha_{8k+2},\ldots,\alpha_{8k+6}]
           (S^{(1)},S^{(2)},(x,H_2,H_4)).
      \end{aligned}
    \]
    Then $\bar Q_{8k+7}$ is decodable given $(H_2,H_4)$.
\end{lemma}
\begin{proof}
    % The input pair contains the unknown coefficients of H_8, so applying the
    % general fill lemma here would be circular.  We instead give the finite
    % top-down decoder explicitly.
    Put $N=8k$ and write
    \[
      \begin{aligned}
      H_2&=x^2+r_1x+r_0,&
      H_4&=x^4+s_3x^3+s_2x^2+s_1x+s_0,\\
      u&=\alpha_1,&v&=\alpha_2,&w&=\alpha_0,&\rho&=\alpha_3.
      \end{aligned}
    \]
    If $k=1$, the defining base of the $T$ recursion gives
    $S^{(1)}=H_8$, $S^{(2)}=\widetilde H_8$, and both internal remainders
    vanish.  If $k\ge2$, then $N<8k+7$, so the admissibility hypothesis
    implies $N$-admissibility and
    \Cref{lem:Rk2l,lem:Rk2l-top-boundary} apply to the internal $T_{k,8}$
    call.  Its level is $l=3$, so it does not use the exceptional shared odd
    base.
    For the outer $A_4$ call put
    \[
      b_i=\alpha_{N+6-i}\ (0\le i\le4),\qquad
      a_i=\alpha_{N-4+i}\ (0\le i\le5).
    \]
    Thus the six $a_i$'s are the consecutive parameters in the $Q_3$ and
    additive slots, while the five $b_i$'s are the scalar slots.  More
    explicitly, the outer expression is
    \[
    \begin{aligned}
      U_0&=(H_4+b_3)S^{(1)}+(x+a_5)(H_2+a_4)+a_3,\\
      V_0&=(H_4+b_4)S^{(2)}+a_2,\\
      C_1&=(H_2+b_1)U_0+a_1,
      &C_2&=(H_2+b_2)V_0+a_0,\\
      \bar Q_{N+7}&=(x+b_0)C_1+C_2.
    \end{aligned}
    \]
    (The notation in this display is just a relabelling of the eleven
    parameters $\alpha_{N-4},\ldots,\alpha_{N+6}$ used by $A_4$.)

    We use the following top-coefficient notation.  If $X$ has degree $N$, let
    \[
      \widehat X(t)=x^{-N}X(x)\big|_{x=t^{-1}},
      \qquad
      p_r=[x^{N+7-r}]\bar Q_{N+7}\quad(0\le r\le7).
    \]
    Equality modulo $t^8$ below means equality of the top eight coefficients.
    For $k=1$ the two $S$-polynomials are exactly $H_8$ and
    $\widetilde H_8$; for $k\ge2$ their remainders have degree $N-8$ by
    \Cref{lem:Rk2l}.  In either case, the coefficients of
    $\widehat S^{(1)}$ and $\widehat S^{(2)}$ through $t^7$ are those of
    $H_8^k$ and $\widetilde H_8^k$, respectively.

    The first three rows have unit pivots:
    \[
      \begin{array}{c|ccc}
        \text{coefficient of }\bar Q_{N+7}&[x^{N+6}]&[x^{N+5}]&[x^{N+4}]\\ \hline
        \text{new parameter}&b_0&b_1&b_2\\
        \text{slope}&1&1&1
      \end{array}
    \]
    Thus, for $i=0,1,2$ in order, set $b_i=p_{i+1}-K_i$, where $K_i$
    is row $p_{i+1}$ of the displayed outer expression with the already
    recovered $b_0,\ldots,b_{i-1}$ substituted and all remaining unknown
    parameters set to zero.  Expanding the monic factors shows that this
    row is exactly $b_i+K_i$: no later parameter reaches it.

    We next solve simultaneously for $(b_3,b_4,u,v)$.  For a compact
    verification of the four pivots, put
    \[
      c_j\coloneqq [x^{N-j}]
        \bigl((H_4+x)(H_4+H_2)\bigr)^k
        \quad(0\le j\le3).
    \]
    These are known from $H_2,H_4$ (and $c_1=2ks_3$).  The coefficients of
    $S^{(1)}$ and $S^{(2)}$ in the first eight degrees below their leading
    term are those of the corresponding powers.  In the normalized variable
    $t=x^{-1}$, the second branch carries an additional factor $t$ because
    it has degree $N+6$ rather than $N+7$; retaining this shift is essential
    in the convolution below.  Since $u$ and $v$ first
    enter $H_8$ four degrees below its leading term, the four rows
    $p_4,p_5,p_6,p_7$ are affine in $(b_3,b_4,u,v)$, where
    $p_r=[x^{N+7-r}]\bar Q_{N+7}$.  After the first three pivots, direct
    convolution gives
    \[
      \begin{pmatrix}p_4\\p_5\\p_6\\p_7\end{pmatrix}
       =\begin{pmatrix}p^0_4\\p^0_5\\p^0_6\\p^0_7\end{pmatrix}
        +M\begin{pmatrix}b_3\\b_4\\u\\v\end{pmatrix},
    \]
    where $p^0_r$ is the coefficient obtained by setting
    $b_3=b_4=u=v=0$ in the displayed outer expression.  These four baseline
    coefficients depend only on $H_2,H_4$ and the already recovered
    $b_0,b_1,b_2$.  The matrix entries are
    \[
    \begin{gathered}
      A_1=b_0+r_1+c_1,\qquad D=r_1+c_1,\\
      C=b_0(c_1+r_1)+c_2+r_1c_1+r_0+b_1,\qquad
      F=c_2+r_1c_1+r_0+b_2,\\
      E=b_0(c_2+r_1c_1+r_0+b_1)+c_3+r_1c_2
        +(r_0+b_1)c_1,\\
      L=b_0+2r_1+(2k-1)s_3+1,
    \end{gathered}
    \]
    and
    \[
      M=
      \begin{pmatrix}
        1&0&k&k\\
        A_1&1&k(A_1+1)&k(A_1+1)\\
        C&D&k(C+D)&k(C+D-1)\\
        E&F&k(E+F-1)&k(E+F-L)
      \end{pmatrix}.
    \]
    Replacing the last two columns by
    $C_3-kC_1-kC_2$ and $C_4-kC_1-kC_2$ gives
    $(0,0,0,-k)^{\mathsf T}$ and $(0,0,-k,-kL)^{\mathsf T}$.
    Since the upper $2\times2$ block in the first two columns has determinant
    one, this proves $\det M=-k^2$.  More importantly, the same reduction
    gives the following explicit inverse.  Put $\Delta p_r=p_r-p^0_r$ and
    successively compute
    \[
      \begin{aligned}
        y_1&=\Delta p_4,&
        y_2&=\Delta p_5-A_1y_1,\\
        v&=(Cy_1+Dy_2-\Delta p_6)/k,&
        u&=(Ey_1+Fy_2-kLv-\Delta p_7)/k,\\
        b_3&=y_1-k(u+v),&
        b_4&=y_2-k(u+v).
      \end{aligned}
    \]
    Indeed the four forward equations become
    \[
      \begin{aligned}
        \Delta p_4&=y_1,&
        \Delta p_5&=A_1y_1+y_2,\\
        \Delta p_6&=Cy_1+Dy_2-kv,&
        \Delta p_7&=Ey_1+Fy_2-ku-kLv,
      \end{aligned}
    \]
    when $y_1=b_3+k(u+v)$ and $y_2=b_4+k(u+v)$.
    Substitution therefore recovers the four input parameters from their
    rows, and sends any four proposed rows back to themselves.  The only
    divisions are by the fixed field constant $k$, which is
    invertible under $(8k+7)$-admissibility; there is no data-dependent
    division.

    It remains to recover $w$ and $\rho$, and then, when $k\ge2$, the
    internal $T$ block.  For $k\ge2$, the next two coefficients include the
    parameter-free boundary rows of the two internal remainders; these are
    known by \Cref{lem:Rk2l-leading-coeff,lem:Rk2l-top-boundary}.  For $k=1$
    the remainders vanish.  Thus in both cases every contribution other than
    the indicated scalar shifts is already known after the preceding steps.
    Indeed, writing
    $H_8=H_8^0+w$, the first term of $(H_8^0+w)^k$ containing $w$ is
    $kw(H_8^0)^{k-1}$, of degree $N-8$; multiplication by the monic
    degree-$7$ factor $A$ places it in degree $N-1$.  Terms containing two
    copies of $w$ are at least eight degrees lower.  Likewise the first term
    containing $\rho$ in $(H_8+\rho)^k$ is
    $k\rho H_8^{k-1}$, and the monic degree-$6$ factor $B$ places it in
    degree $N-2$.  Therefore
    \[
      [x^{N-1}]\bar Q_{N+7}=kw+K_{N-1},\qquad
      [x^{N-2}]\bar Q_{N+7}=k\rho+K_{N-2}(w),
    \]
    for known polynomials $K_{N-1},K_{N-2}$, computed by setting the current
    scalar to zero and using the supplied remainder boundary coefficients.
    The explicit pivots are
    \[
      w=\bigl([x^{N-1}]\bar Q_{N+7}-K_{N-1}\bigr)/k,
      \qquad
      \rho=\bigl([x^{N-2}]\bar Q_{N+7}-K_{N-2}(w)\bigr)/k.
    \]
    The first boundary coefficient of $H_8^k$ contributes $kw$; the
    second branch contributes $k(w+\rho)$ one row later.
    Consequently $H_8$ and $\widetilde H_8$ are now known.

    Subtract the known main terms and set
    \[
      A=(x+b_0)(H_2+b_1)(H_4+b_3),\qquad
      B=(H_2+b_2)(H_4+b_4).
    \]
    Then
    \[
      \bar Q_{N+7}-AH_8^k-B\widetilde H_8^k
       =AR^{(1)}_{k,8}+BR^{(2)}_{k,8}+F,\qquad \deg F\le6,
    \]
    where
    \[
      F=(x+b_0)\bigl((H_2+b_1)Q_3+a_1\bigr)+(H_2+b_2)a_2+a_0.
    \]
    If $k=1$, then $R^{(1)}_{1,8}=R^{(2)}_{1,8}=0$ and the internal parameter
    interval $\alpha_4,\ldots,\alpha_{N-5}$ is empty, so the displayed
    residual is simply $F$.  Suppose now that $k\ge2$.  Since $A$ and $B$
    are monic of degrees $7$ and $6$, respectively, the
    coefficient in degree $j+6$ of the first two terms is
    \[
      [x^j]\bigl(xR^{(1)}_{k,8}+R^{(2)}_{k,8}\bigr)
      +\text{a polynomial in higher coefficients of }R^{(1)},R^{(2)}.
    \]
    Descending through degrees $N-3,N-4,\ldots,6$ therefore transfers the
    triangular decoder of \Cref{lem:Rk2l}(3) to the internal parameter block
    $\alpha_4,\ldots,\alpha_{N-5}$.  This explicitly handles the overlap of
    the two multiplied remainder polynomials; no fill-compatibility lemma is
    invoked.  At the final row (degree $6$, corresponding to the inner
    coefficient of degree $0$), the low term has only its known leading
    contribution $[x^6]F=1$; subtracting this constant leaves the same
    triangular pivot.  The unknown coefficients of $F$ start in degree at
    most $5$.

    (For $k=1$ this entire internal descent is void.)  After reconstructing
    $S^{(1)},S^{(2)}$, compute the low residual $F$.
    Its coefficients in degrees $5,4,3,2,1,0$ have the unit-triangular pivot
    table
    \[
      \begin{array}{c|cccccc}
        \text{degree}&5&4&3&2&1&0\\ \hline
        \text{new parameter}&a_5&a_4&a_3&a_2&a_1&a_0\\
        \text{slope}&1&1&1&1&1&1.
      \end{array}
    \]
    The coefficient in degree $6$ is the known monic leading term.  An
    explicit inverse for these six rows is monic division by
    $J=(x+b_0)(H_2+b_1)$: the quotient is $Q_3$, since
    \[
      F=JQ_3+(x+b_0)a_1+(H_2+b_2)a_2+a_0
    \]
    and the last three terms have degree at most $2<\deg J$.
    Apply the explicit $Q_3$ decoder to the quotient to recover
    $a_5,a_4,a_3$.  For the remainder $R=F-JQ_3$, compute
    \[
      a_2=[x^2]R,\qquad a_1=[x]R-r_1a_2,\qquad
      a_0=[x^0]R-b_0a_1-(r_0+b_2)a_2.
    \]
    Monic division and these assignments give both inverse compositions
    directly and recover the six parameters in descending order.
    Together with the already
    recovered internal block and $(w,u,v,\rho)$ this is every parameter in
    $\bar Q_{N+7}$, proving decodability given $(H_2,H_4)$.
\end{proof}

\begin{corollary}[The degree-$15$ barred gadget]\label{lem:barQ15}
    The polynomial $\bar Q_{15}$ is decodable given $(H_2,H_4)$ over every
    $15$-admissible field.
\end{corollary}
\begin{proof}
    Specialize \Cref{lem:barQ8k+7} to $k=1$.
\end{proof}

\begin{lemma}[Odd-degree gadgets from $(H_2,H_4)$]\label{lem:odd-gadgets-H2H4}
    Let $H_2,H_4\in\mathbb{F}[x]$ be monic with $\deg H_2=2$ and $\deg H_4=4$.
    For each degree $n$ asserted below, assume that $\mathbb F$ is
    $n$-admissible.
    Then:
    \begin{enumerate}
        \item If $n\equiv 1\pmod 4$ and $n\ge 5$, then $Q_n$ is decodable given $H_2$ by \Cref{lem:Q4k+1-from-H2}.
        \item The base degrees $n=3$ and $n=7$ are the explicit $Q_3$ and
        $Q_7$ constructions above.  If
        $n\equiv 3\pmod 8$ and $n\ge 11$, then $Q_n$ is decodable given
        $(H_2,H_4)$ by the known-powers construction with $l=2$ (i.e.
        \Cref{alg:constr-Q-odd} and \Cref{lem:Q-odd-degree-with-powers}
        specialized to $l=2$).
        \item If $n\equiv 7\pmod 8$ and $n\ge 15$, then $\bar Q_n$ is
        decodable given $(H_2,H_4)$ by \Cref{lem:barQ8k+7}; its $k=1$
        specialization is the degree-$15$ gadget
        \Cref{lem:barQ15}.
    \end{enumerate}
\end{lemma}
\begin{proof}
    For $n\equiv 1\pmod 4$, write $n=4k+1$ and apply \Cref{lem:Q4k+1-from-H2}.

    The cases $n=3,7$ are the explicit $Q_3,Q_7$ constructions.  For
    $n\equiv 3\pmod 8$
    with $n\ge11$, write $n=2^{l+1}k+(2^l-1)$ with $l=2$ and apply the
    construction \Cref{alg:constr-Q-odd} together with
    \Cref{lem:Q-odd-degree-with-powers} (specialized to $l=2$).

    For $n\equiv 7\pmod 8$ and $n\ge 15$, write $n=8k+7$ with $k\ge1$
    and apply \Cref{lem:barQ8k+7}.
\end{proof}

\paragraph{Selected odd gadget notation.}
In the two final induction steps, the degree of an auxiliary odd gadget is
not always of the Mersenne form for which the required higher powers have
already been built.  To avoid hiding this distinction in the algorithm
captions, write $\mathcal Q_d$ for the following degree-$d$ gadget (with a
fresh, consecutive parameter block):
\[
\mathcal Q_d\coloneqq
\begin{cases}
 Q_1, & d=1,\\
 Q_3, & d=3,\\
 Q_7, & d=7,\\
 Q_d, & d\equiv1\pmod4,\\
 Q_d, & d\equiv3\pmod8,\\
 \bar Q_d, & d\equiv7\pmod8\text{ and }d\ge15.
\end{cases}
\]
Here the first two non-base lines use respectively
\Cref{lem:Q4k+1-from-H2} and the known-powers construction with
$(H_2,H_4)$; the first three lines are the explicit $Q_1,Q_3,Q_7$ bases,
and the final line uses the barred construction
\Cref{lem:barQ15,lem:barQ8k+7}.  If additional powers are available, the
standard known-powers $Q_d$ may be used instead of $\mathcal Q_d$.
In every case $\mathcal Q_d$ is monic, decodable from the displayed known
powers, and has the cost $\lfloor d/2\rfloor$ recorded below.

\begin{lemma}[Exact costs of the auxiliary odd gadgets]
\label{lem:odd-gadgets-count}
    Assume the indicated known powers are already available.
    \begin{enumerate}
      \item $\mathcal Q_{4k+1}(x,H_2)=Q_{4k+1}(x,H_2)$ uses $2k$
      multiplications.
      \item For $l\ge2$, the known-powers construction
      $Q_{2^{l+1}k+(2^l-1)}$ uses
      \[
        2^l k+2^{l-1}-1
        =\left\lfloor\frac{2^{l+1}k+(2^l-1)}2\right\rfloor
      \]
      multiplications.
      \item $\mathcal Q_{8k+7}(x,H_2,H_4)=\bar Q_{8k+7}(x,H_2,H_4)$ uses
      $4k+3$ multiplications; this includes $\bar Q_{15}$ when $k=1$.
    \end{enumerate}
    Consequently every auxiliary odd-degree $Q$-gadget used in the two final
    induction steps has cost $\lfloor d/2\rfloor$, where $d$ is its degree.
\end{lemma}
\begin{proof}
    The first construction consists of the call $T_{2k,2}$, whose cost is
    $2k-1$ by \Cref{lem:T-multiplication-count}, followed by the one product
    $(x+\beta)T^{(1)}$.

    For the second construction, forming
    $\hat H_{2^l}=H_{2^l}+Q_{2^{l-1}-1}$ costs
    $2^{l-2}-1$ multiplications.  The internal $T_{2k,2^l}$ call costs
    $(2k-1)2^{l-1}$, and the outer $A_{2^{l-1}}$ call costs
    $3\cdot2^{l-2}$ (the formula in \Cref{lem:fill-Q-count}, with the direct
    $A_2$ base when $l=2$).  Their sum is
    \[
      (2^{l-2}-1)+(2k-1)2^{l-1}+3\cdot2^{l-2}
      =2^lk+2^{l-1}-1.
    \]

    Finally, $\bar Q_{8k+7}$ uses one product to form $H_8$, then
    $4(k-1)$ products for $T_{k,8}$, and six products for $A_4$.  Thus its
    cost is $1+4(k-1)+6=4k+3$.  The remaining auxiliary cases
    $Q_{2^r-1}$ have cost $2^{r-1}-1=\lfloor(2^r-1)/2\rfloor$ by
    \Cref{lem:fill-Q-count}; together these cases cover the gadgets listed in
    \Cref{lem:odd-gadgets-H2H4} (including $Q_1,Q_3,Q_7$).
\end{proof}

\subsubsection{The \texorpdfstring{$8k+7$}{8k+7} step}

\Cref{fig:odd-induction-steps} pictures the two induction steps.

\begin{figure}[H]
  \centering
  \resizebox{\textwidth}{!}{% Nested-shell view of the 8k+7 and 8k+3 induction steps.
\begin{tikzpicture}[x=1cm,y=1cm,>=Latex]
  \draw[fp decode] (0,5.2) -- node[fp label, above] {descending coefficient degree}
    (13.55,5.2);
  \node[fp label, anchor=east] at (0,5.2) {high};
  \node[fp label, anchor=west] at (13.55,5.2) {low};

  \node[fp panel title] at (0,4.67) {(a) The $8k+7$ step: two nested square shells};
  \node[fp active, minimum width=49mm, anchor=west] (s3) at (0,3.92)
    {$xS_3^2+(S_3+b)^2$\\recover $S_3,b$};
  \node[fp seam, minimum width=8mm, anchor=west] at (4.9,3.92) {$-1$};
  \node[fp second, minimum width=36mm, anchor=west] (s2) at (5.7,3.92)
    {$xS_2^2+(S_2+a)^2$\\negate residual\\recover $S_2,a$};
  \node[fp seam, minimum width=8mm, anchor=west] at (9.3,3.92) {$-1$};
  \node[fp recursive, minimum width=30mm, anchor=west] (small7) at (10.1,3.92)
    {$-P_{2k+1}$\\negate; recurse};
  \draw[fp dependence] (s3.south east) to[bend right=10] (s2.south west);
  \draw[fp dependence] (s2.south east) to[bend right=10] (small7.south west);
  \node[fp panel title] at (0,2.95) {(b) The $8k+3$ step: shell, squared recursion, delayed low gadget};
  \node[fp active, minimum width=45mm, anchor=west] (top3) at (0,2.2)
    {$xS_2^2+(S_2+a)^2$\\recover $S_2,a$};
  \node[fp seam, minimum width=8mm, anchor=west] at (4.5,2.2) {$-1$};
  \node[fp recursive, minimum width=45mm, anchor=west] (sqrec) at (5.3,2.2)
    {$-x(S^{(1)}_1)^2-(S^{(2)}_1)^2$\\recover squares; take monic roots};
  \node[fp second, minimum width=30mm, anchor=west] (low3) at (9.8,2.2)
    {$xS_3+\alpha_0$\\low residual};
  \draw[fp dependence] (top3.south east) to[bend right=10] (sqrec.south west);
  \draw[fp dependence] (sqrec.south east) to[bend right=10] (low3.south west);
  \node[fp known, minimum width=37mm] (powers) at (7.4,0.78)
    {$P_{2k+1}$ decoder: $H_2$\\then $S_2$ decoder: $H_4$};
  \draw[fp decode] (sqrec.south) -- (powers.north);
  \draw[fp dependence] (powers.east) -- (low3.south west);
  \node[fp label, anchor=west, text=fpBlue!80!black] at (0,0.02)
    {Both steps recover auxiliary polynomials before invoking their conditional parameter decoders.};
\end{tikzpicture}
}
  \caption{The two recursive odd-degree steps.  For $8k+7$, two square
  shells are peeled in succession and leave the smaller polynomial
  $P_{2k+1}$.  For $8k+3$, the outer shell exposes a compatible squared
  smaller pair; its monic square roots recover the smaller instance, whose
  decoder reconstructs $H_2$.  Decoding $S_2=Q_{4k+1}(x,H_2)$ next
  reconstructs $H_4$, after which the low auxiliary gadget can be decoded.
  Red cells mark the square-gadget boundary corrections; the second shell
  in the $8k+7$ panel is read after negating the first residual.  The
  additional row-$2k$ correction for the squared recursion in the $8k+3$
  step is given in the proof.}
  \label{fig:odd-induction-steps}
\end{figure}

\begin{algorithm}[H]
    \caption{
        Extending a compatible joint realization from $2k+1$ to $8k+7$.
        The odd-degree auxiliary calls below use the selected gadgets
        $\mathcal Q_d$ of \Cref{lem:odd-gadgets-H2H4}; they require only the
        displayed $H_2,H_4$ (and only $H_2$ in the $d\equiv1\pmod4$ case).
    }
    \label{alg:constr-8k+7}

    \begin{algorithmic}
        \State $S^{(1)}_1 = T^{(1)}_{2k + 1}[\alpha_{0}, \ldots, \alpha_{2k}](x)$ \Comment{The smaller construction records the $H_2,H_4$ used below.}
        \State $S^{(1)}_2 = \mathcal Q_{2k + 1}[\alpha_{2k + 1}, \ldots, \alpha_{4k + 1}](x, H_2(x), H_4(x))$
        \State $S^{(1)}_3 = \mathcal Q_{4k + 3}[\alpha_{4k + 2}, \ldots, \alpha_{8k + 4}](x, H_2(x), H_4(x))$
        \State Put $s=\alpha_{8k+5}$, $d=\alpha_{8k+6}$ and, for the
        decoder, $a=(s-d)/2$, $b=(s+d)/2$.
        \State $T^{(1)}_{8k + 7}[\alpha_0, \ldots, \alpha_{8k + 6}](x)
            = (S^{(1)}_3 + S^{(1)}_2)(S^{(1)}_3 - S^{(1)}_2) + S^{(1)}_1
            = (S^{(1)}_3)^2 - (S^{(1)}_2)^2 + S^{(1)}_1$
        \\
        \State $S^{(2)}_1 = T^{(2)}_{2k + 1}[\alpha_{0}, \ldots, \alpha_{2k}](x)$
        \State For the identities put $S^{(2)}_2=S^{(1)}_2+a$ and
        $S^{(2)}_3=S^{(1)}_3+b$; the circuit does not materialize these two
        shifted polynomials.
        \State $T^{(2)}_{8k + 7}[\alpha_0, \ldots, \alpha_{8k + 6}](x)
            = \bigl((S^{(1)}_3+S^{(1)}_2)+s\bigr)
              \bigl((S^{(1)}_3-S^{(1)}_2)+d\bigr)+S^{(2)}_1
            = (S^{(2)}_3 + S^{(2)}_2)(S^{(2)}_3 - S^{(2)}_2) + S^{(2)}_1
            = (S^{(2)}_3)^2 - (S^{(2)}_2)^2 + S^{(2)}_1$

        \State \begin{align}
            P_{8k + 7}[\alpha_0, \ldots, \alpha_{8k + 6}](x)
            &= x T^{(1)}_{8k + 7} + T^{(2)}_{8k + 7}
            \\&=x(S^{(1)}_3)^2+(S^{(1)}_3+b)^2
                 -x(S^{(1)}_2)^2\notag\\
            &\quad -(S^{(1)}_2+a)^2
                 +xS^{(1)}_1+S^{(2)}_1
        \end{align}
    \end{algorithmic}
\end{algorithm}

\begin{lemma}[The $8k+7$ induction step is decodable (and preserves compatibility)]\label{lem:8k+7-splittable}
    Let $k\ge2$ and assume that $\mathbb{F}$ is $(8k+7)$-admissible.
    Suppose that $P_{2k+1}=xT^{(1)}_{2k+1}+T^{(2)}_{2k+1}$ is decodable, and that the auxiliary polynomials $S^{(1)}_2,S^{(1)}_3$ in \Cref{alg:constr-8k+7} are decodable given the relevant known-power data.
    Assume moreover that the underlying pair
    $(T^{(1)}_{2k+1},T^{(2)}_{2k+1})$ is compatible on some window $G$
    with no auxiliary data, and that its construction records the monic
    polynomials $H_2,H_4$ used by the two auxiliary gadgets as polynomial
    byproducts of the smaller parameter block.
    Define
    \[
      G_{8k+7}\coloneqq \{2k+1,\ldots,4k+3\}\cup\{4k+3,\ldots,8k+6\}\cup G.
    \]
    Then the polynomial $P_{8k+7}$ produced by the construction is decodable,
    and the output pair $(T^{(1)}_{8k+7},T^{(2)}_{8k+7})$ is compatible on
    $G_{8k+7}$ with no auxiliary data.  If the smaller pair has a joint
    $m$-realization recording $H_2,H_4$, then the displayed construction has
    a joint $(m+3k+3)$-realization and records the same two powers.
\end{lemma}
\begin{proof}
    Write $S_2\coloneqq S^{(1)}_2$ and
    $S_3\coloneqq S^{(1)}_3$ for brevity.  The two scalar parameter
    coordinates are $s=\alpha_{8k+5}$ and $d=\alpha_{8k+6}$; put
    $a=(s-d)/2$ and $b=(s+d)/2$, so that $a+b=s$ and $b-a=d$.
    This is an invertible fixed linear change of coordinates because the
    admissibility hypothesis makes $2$ invertible.

    Since $S_3=\mathcal Q_{4k+3}$ is monic of degree $4k+3$, the polynomial $xS_3^2+(S_3+b)^2$ has degree $8k+7$.
    All remaining summands in the explicit formula for $P_{8k+7}$ have degree at most $4k+3$:
    indeed, $S_2=\mathcal Q_{2k+1}$ has degree $2k+1$, so $xS_2^2$ has degree $4k+3$ and $(S_2+a)^2$ has degree $4k+2$, while $xS^{(1)}_1+S^{(2)}_1=P_{2k+1}$ has degree $2k+1$.
    Therefore we may write
    \[
        P_{8k+7} = \underbrace{xS_3^2 + (S_3+b)^2}_{\text{square gadget at degree }4k+3}
            \;+\; E_3
    \]
    where $\deg E_3\le 4k+3$.
    The boundary coefficient of this error is fixed:
    $[x^{4k+3}]E_3=-1$.  Therefore
    \Cref{lem:square-gadget-boundary}, with this supplied boundary
    coefficient, derives $S_3$ and $b$ from $P_{8k+7}$.

    Subtracting the derived square-gadget polynomial $xS_3^2+(S_3+b)^2$ from $P_{8k+7}$ leaves
    \[
        P' \coloneqq -xS_2^2 - (S_2+a)^2 + P_{2k+1}.
    \]
    Equivalently, $-P' = xS_2^2 + (S_2+a)^2 + E_2$ with $E_2=-P_{2k+1}$ and $\deg E_2\le 2k+1=\deg S_2$.
    Its boundary coefficient is again fixed, $[x^{2k+1}]E_2=-1$.
    Applying \Cref{lem:square-gadget-boundary}, with this boundary
    coefficient, derives $S_2$ and $a$ from $P'$, and by
    construction both polynomials are now available without using their
    internal known-power data.

    Finally,
    \[
        P_{2k+1} = P' + xS_2^2 + (S_2+a)^2
    \]
    is derivable from $P_{8k+7}$, so decodability of $P_{2k+1}$ extracts all
    parameters of the $2k+1$ block.  We can then recompute the recorded
    byproducts $H_2,H_4$ of that smaller construction.  Only at this point do
    we invoke decodability of the auxiliary gadgets: from the already derived
    polynomials $S_2,S_3$ and their now-known powers,
    \Cref{lem:discharge-side-information} extracts their parameter blocks;
    the reconstructed $H_2,H_4$ are thereby discharged rather than silently
    treated as fixed auxiliary data.
    Together with $s=a+b$ and $d=b-a$, this extracts
    $\alpha_0,\ldots,\alpha_{8k+6}$ and proves decodability.

    For compatibility we retain the row cutoffs in those two peelings.  If
    $S$ is monic of degree $e$, the descending recurrence in
    \Cref{lem:square-gadget-boundary} first reads the coefficient
    $[x^q]S$ in row $e+q+1$; the scalar shift is read in row $e$.
    Consequently $[x^j]S^2$ is recovered using only rows at least $j+1$,
    while $[x^j](S+\delta)^2$ uses only rows at least $j$.  Subtracting a
    recovered square gadget at row $j$ also uses only rows at least $j$, so
    the same statement remains true for the second peeling and for the
    residual $P_{2k+1}$.

    Now apply the unconditional compatibility hypothesis for the smaller
    pair to that residual.  It recovers
    $[x^j]T^{(1)}_{2k+1}$ from output rows at least $j+1$ and
    $[x^j]T^{(2)}_{2k+1}$ from rows at least $j$.  Substitution in
    \[
      T^{(1)}_{8k+7}=S_3^2-S_2^2+T^{(1)}_{2k+1},\qquad
      T^{(2)}_{8k+7}=(S_3+b)^2-(S_2+a)^2+T^{(2)}_{2k+1}
    \]
    gives exactly the two required cutoffs on $G_{8k+7}$.  No coefficients
    of the internal powers used to define $S_2,S_3$ occur in this
    reconstruction, so the output compatibility is unconditional.

    Finally run the smaller $m$-realization once.  The two selected auxiliary
    gadgets cost $k$ and $2k+1$ products by
    \Cref{lem:odd-gadgets-count}; they reuse the recorded $H_2,H_4$ wires.
    The two displayed difference-of-squares outputs cost one product each.
    Hence the joint cost is $m+k+(2k+1)+2=m+3k+3$, and the inherited power
    wires remain recorded outputs.
\end{proof}

\begin{example}[Worked example: $k=7$ (constructing $63$)]
    Instantiating the construction with $k=7$ gives $2k+1=15$, $4k+3=31$, and $8k+7=63$.
    Let $(T^{(1)}_{15},T^{(2)}_{15})$ be the compatible seven-product joint
    realization from \Cref{sec:special-cases}, which records $H_2,H_4$, and
    let $\bar Q_{15}$ and
    $\bar Q_{31}$ be the selected degree-$15$ and degree-$31$ gadgets from
    the convention above.
    Put $s\coloneqq\alpha_{61}$, $d\coloneqq\alpha_{62}$ and
    $a=(s-d)/2$, $b=(s+d)/2$.  The construction specializes to
    \[
      \begin{aligned}
        T^{(1)}_{63}
          &=\bar Q_{31}^2-\bar Q_{15}^2+T^{(1)}_{15},\\
        T^{(2)}_{63}
          &=(\bar Q_{31}+\bar Q_{15}+s)
            (\bar Q_{31}-\bar Q_{15}+d)+T^{(2)}_{15}\\
          &=(\bar Q_{31}+b)^2-(\bar Q_{15}+a)^2+T^{(2)}_{15}.
      \end{aligned}
    \]
    and hence
    \[
        P_{63}=xT^{(1)}_{63}+T^{(2)}_{63}
        = x\bar Q_{31}^2 + (\bar Q_{31}+b)^2
          -x\bar Q_{15}^2-(\bar Q_{15}+a)^2+P_{15}.
    \]
    The parameter blocks are
    \[
      \begin{array}{c|c}
        \alpha_0,\ldots,\alpha_{14}&(T^{(1)}_{15},T^{(2)}_{15})\\
        \alpha_{15},\ldots,\alpha_{29}&\bar Q_{15}\\
        \alpha_{30},\ldots,\alpha_{60}&\bar Q_{31}\\
        \alpha_{61},\alpha_{62}&s,d.
      \end{array}
    \]
\end{example}

\subsubsection{The \texorpdfstring{$8k+3$}{8k+3} step}
\begin{algorithm}[H]
    \caption{
        Extending a compatible joint realization from $2k+1$ to $8k+3$.
        The call to $\mathcal Q_{2k-1}$ below uses the selected gadget of
        \Cref{lem:odd-gadgets-H2H4}: it needs only the $H_2,H_4$ made
        derivable by the preceding two calls (with the explicit $Q_3,Q_7$
        bases in the small cases).
    }
    \label{alg:constr-8k+3}

    \begin{algorithmic}
        \State $S^{(1)}_1 = T^{(1)}_{2k + 1}[\alpha_{2k}, \ldots, \alpha_{4k}](x)$ \Comment{The smaller construction records $H_2(x)$.}
        \State $S^{(1)}_2 = Q_{4k + 1}[\alpha_{4k + 2}, \ldots, \alpha_{8k + 2}](x, H_2(x))$ \Comment{This call records the quartic $H_4(x)$.}
        \If{$k > 1$}
            \State $S^{(1)}_3 = \mathcal Q_{2k - 1}[\alpha_{1}, \ldots, \alpha_{2k - 1}](x, H_2(x), H_4(x))$
        \Else
            \State $S^{(1)}_3 = \alpha_1$
        \EndIf
        \State $T^{(1)}_{8k + 3}[\alpha_0, \ldots, \alpha_{8k + 2}](x)
            = (S^{(1)}_2 + S^{(1)}_1)(S^{(1)}_2 - S^{(1)}_1) + S^{(1)}_3
            = (S^{(1)}_2)^2 - (S^{(1)}_1)^2 + S^{(1)}_3$
        \\
        \State $S^{(2)}_1 = T^{(2)}_{2k + 1}[\alpha_{2k}, \ldots, \alpha_{4k}](x)$
        \State $S^{(2)}_2 = S^{(1)}_2 + \alpha_{4k + 1}$
        \State $S^{(2)}_3 = \alpha_{0}$
        \State $T^{(2)}_{8k + 3}[\alpha_0, \ldots, \alpha_{8k + 2}](x)
            = (S^{(2)}_2 + S^{(2)}_1)(S^{(2)}_2 - S^{(2)}_1) + S^{(2)}_3
            = (S^{(2)}_2)^2 - (S^{(2)}_1)^2 + S^{(2)}_3$

        \State \begin{align}
            P_{8k + 3}[\alpha_0, \ldots, \alpha_{8k + 2}](x)
            &= x T^{(1)}_{8k + 3} + T^{(2)}_{8k + 3}
            \\&=x(S^{(1)}_2)^2+(S^{(1)}_2+\alpha_{4k+1})^2
                 -x(S^{(1)}_1)^2\notag\\
            &\quad -(S^{(2)}_1)^2+xS^{(1)}_3+S^{(2)}_3
        \end{align}
    \end{algorithmic}
\end{algorithm}

\begin{lemma}[The $8k+3$ induction step is decodable (and preserves compatibility)]\label{lem:8k+3-splittable}
    Let $k\ge1$ and assume that $\mathbb{F}$ is $(8k+3)$-admissible.
    Suppose that $P_{2k+1}=xT^{(1)}_{2k+1}+T^{(2)}_{2k+1}$ is decodable,
    and that the underlying pair
    $(T^{(1)}_{2k+1},T^{(2)}_{2k+1})$ is compatible on some $G$ with no
    auxiliary data.  Assume that the smaller construction records its monic
    quadratic $H_2$ as a polynomial byproduct.  Assume also that the auxiliary polynomials
    $S^{(1)}_2$ and $S^{(1)}_3$ in \Cref{alg:constr-8k+3} are decodable once
    their indicated known powers are supplied.
    Define
    \[
      G_{8k+3}\coloneqq
      \{4k+1,\ldots,8k+2\}\cup (2k+G)\cup \rng{2k+1}.
    \]
    Then the polynomial $P_{8k+3}$ produced by the construction is decodable,
    and the output pair $(T^{(1)}_{8k+3},T^{(2)}_{8k+3})$ is compatible on
    $G_{8k+3}$ with no auxiliary data.  If the smaller pair has a joint
    $m$-realization recording $H_2$, then the displayed construction has a
    joint $(m+3k+1)$-realization and records both $H_2$ and the quartic $H_4$
    formed by its $Q_{4k+1}$ call.
\end{lemma}
\begin{proof}
    Write $S_2\coloneqq S^{(1)}_2$, $S_3\coloneqq S^{(1)}_3$, and $a\coloneqq \alpha_{4k+1}$ for the scalar shift.

    From the explicit formula in the algorithm we have
    \[
        P_{8k+3} = \underbrace{xS_2^2 + (S_2+a)^2}_{\text{square gadget at degree }4k+1} + E,
    \]
    where $\deg E\le 4k+1$ (all remaining terms involve only the $2k+1$-block and $xS_3+\alpha_0$).
    The boundary coefficient of this error is fixed:
    $[x^{4k+1}]E=-1$.  Therefore
    \Cref{lem:square-gadget-boundary}, with this supplied boundary
    coefficient, derives $S_2$ and $a$ from $P_{8k+3}$.

    Subtracting the derived square-gadget polynomial
    $xS_2^2+(S_2+a)^2$ yields
    \[
      P'=-\Psi+xS_3+\alpha_0,
      \qquad
      \Psi\coloneqq x(S^{(1)}_1)^2+(S^{(2)}_1)^2,
    \]
    where $S^{(i)}_1=T^{(i)}_{2k+1}$.  The last two terms have degree at
    most $2k$.  Hence $[x^d]\Psi=-[x^d]P'$ for $d>2k$.  At $d=2k$ the
    only correction is the fixed leading coefficient of $xS_3$: it is zero
    for $k=1$, and one for $k>1$.  Thus every coefficient of $\Psi$ in the
    window $2k+G$ is recoverable from $P'$ without knowing any coefficient of
    the internal powers.

    The unconditional compatibility hypothesis and
    \Cref{lem:compatible-power} show that
    $((S^{(1)}_1)^2,(S^{(2)}_1)^2)$ is compatible on $2k+G$, again with no
    auxiliary data.  Its decoder therefore recovers both squares from the
    displayed $\Psi$-window.  Taking the descending monic square roots
    recovers $S^{(1)}_1,S^{(2)}_1$, and hence
    $P_{2k+1}=xS^{(1)}_1+S^{(2)}_1$.  Decodability of $P_{2k+1}$ now extracts
    its parameter block, after which the recorded polynomial $H_2$ can be
    recomputed.

    Finally, with $(S^{(1)}_1,S^{(2)}_1)$ derived, we can isolate
    \[
        xS_3+\alpha_0 = P' + x(S^{(1)}_1)^2 + (S^{(2)}_1)^2.
    \]
    The constant term of $xS_3+\alpha_0$ equals $\alpha_0$, so $\alpha_0$ is extractable.
    Subtracting $\alpha_0$ then derives $xS_3$, and hence $S_3$ (by shifting coefficients).
    Decode the already recovered polynomial
    $S_2=Q_{4k+1}(x,H_2)$ first.  By
    \Cref{lem:Q4k+1-from-H2}, this both extracts the parameter block of $S_2$
    and reconstructs its internal quartic $H_4$.  With $(H_2,H_4)$ now
    available, decodability of $S_3$ extracts its parameter block (for $k=1$,
    $S_3=\alpha_1$ is already read directly).  Thus neither auxiliary decoder
    is invoked before the powers it requires have actually been reconstructed;
    \Cref{lem:discharge-side-information} then substitutes those reconstructed
    powers into the two conditional decoders.

    Together with $a$ and $\alpha_0$, these are all parameters, proving
    decodability.

    It remains to verify the causal cutoffs.  For a monic degree-$e$
    polynomial $S$, the square-gadget recurrence first reads $[x^q]S$ in
    row $e+q+1$ and reads its scalar shift in row $e$.  Consequently
    $[x^j]S^2$ is recoverable from output rows at least $j+1$, while
    $[x^j](S+\delta)^2$ is recoverable from rows at least $j$.  This applies
    to $S_2$ on the top window $\{4k+1,\ldots,8k+2\}$.

    For the middle block, recovering $[x^r]\Psi$ from $P'$ uses only the
    output row $r$ and the fixed monic boundary correction at $r=2k$.
    The cutoff conclusion of \Cref{lem:compatible-power} therefore gives
    \[
      \begin{aligned}
      [x^j](S^{(1)}_1)^2
        &\polyfrom([x^r]P_{8k+3}:r\in G_{8k+3},\ r\ge j+1),\\
      [x^j](S^{(2)}_1)^2
        &\polyfrom([x^r]P_{8k+3}:r\in G_{8k+3},\ r\ge j).
      \end{aligned}
    \]
    Finally, after those squares have been subtracted,
    $xS_3+\alpha_0$ is read coefficient by coefficient on
    $\rng{2k+1}$: $[x^j]S_3$ is read in row $j+1$, and $\alpha_0$ in row
    zero.  Substitution in
    \[
      T^{(1)}_{8k+3}=S_2^2-(S^{(1)}_1)^2+S_3,
      \qquad
      T^{(2)}_{8k+3}=(S_2+a)^2-(S^{(2)}_1)^2+\alpha_0
    \]
    gives the required first-component cutoff $j+1$ and second-component
    cutoff $j$.  No internal known power occurs in these formulas, so the
    output compatibility is unconditional.

    For the joint cost, run the smaller circuit once.  The selected gadgets
    of degrees $4k+1$ and $2k-1$ cost $2k$ and $k-1$ products by
    \Cref{lem:odd-gadgets-count} (the second cost is zero when $k=1$), and
    the two displayed difference-of-squares outputs cost one product each.
    The total is $m+2k+(k-1)+2=m+3k+1$.  The first auxiliary call forms and
    records $H_4$, while the smaller circuit's $H_2$ wire is retained.
\end{proof}
